\section{Exact symplectic sewing of annular chains}
\label{sew:section}

This section isolates the analytic construction needed in the geometric
arguments.  The domains are cylinders, not discs.  Consequently, a closed
holomorphic one-form on an overlap need not be exact.  Its period is the
obstruction that must be removed before the iteration can converge.

Put $\mathcal C=\C/\Z$, and write $Y=\operatorname{Im}w$ on
$\mathcal C$.  We use the form
\[
                 \omega_0=dw\wedge dp
\]
on products of $\mathcal C$ with a parameter disc.  A symplectic map
$F=(W,P)$ between such products is called \emph{exact} if
\[
                F^*(P\,dW)-p\,dw
\]
is an exact holomorphic one-form.  Since a cylinder times a disc has first
homology $\Z$, this is equivalent to the vanishing of its integral around
one positively oriented period circle.  All lifts below have degree one:
$W(w+1,p)=W(w,p)+1$ and $P(w+1,p)=P(w,p)$.

\begin{theorem}[Annular-chain sewing]\label{sew:chain}
Let $M$ be a complex surface with a nowhere-zero holomorphic two-form
$\omega$.  Fix $0<h<1/8$, $r>0$, and $A<\infty$.  For each $j\in\Z$,
let $L_j\geq1$ and let
\[
 \phi_j:\{ -2h<Y<L_j+2h\}\times\Delta_r\longrightarrow M
\]
be a holomorphic map with $\phi_j^*\omega=\omega_0$.  Suppose there are
holomorphic translations
\[
 J_j(w,p)=(w+a_j(p),p),\qquad
 |\operatorname{Im}a_j(p)+L_j|<h/8,\qquad |a_j'(p)|\leq A,
\]
and exact symplectic transitions $F_j$ on
\[
 S_j(h,r)=\{|Y-L_j|<h\}\times\Delta_r
\]
such that $\phi_j=\phi_{j+1}\circ F_j$.  Write
\[
                  J_j^{-1}F_j=(w+u_j,p+v_j).
\]
Assume these lifts exist and
$\sup_j\|(u_j,v_j)\|_{S_j(h,r)}\leq e$.

There are $b>0$ and $e_*>0$, depending only on $h,r,A$ and the indicated margins,
such that, if $e<e_*$, these charts give a holomorphic symplectic
immersion
\[
                  \Phi:\C\times\Delta_b\longrightarrow M
\]
which is periodic of period one in the first variable.
On a smaller domain in every original block, with fixed positive
margins at both ends, $\Phi$ is obtained from $\phi_j$ by a small
periodic symplectic change of variables.  The corrected transitions are
\[
        (w,p)\longmapsto(w+\widetilde a_j(p),p).
\]
The changes of variables and $\widetilde a_j-a_j$, together with any
fixed finite number of derivatives on further smaller domains, tend
uniformly to zero as $e\to0$.  None of these assertions requires an
upper bound for the $L_j$ or a bound on the number of blocks.
\end{theorem}

Here and below $\Delta_r=\{p\in\C:|p|<r\}$; a norm on a noncompact
cylinder is the supremum norm of a periodic function.  We give the
estimates and the domain construction explicitly.

\subsection{Domains and the infinitesimal obstruction}

For a current collection of translations, define
\[
\begin{split}
 S_j(s,\rho)&=\{|Y-L_j|<s,\ |p|<\rho\},\\
 D_j(s,\rho)&=\{L_{j-1}-s<Y-\operatorname{Im}a_{j-1}(p),\quad
                         Y<L_j+s,\quad |p|<\rho\}.
\end{split}
\]
The lower edge of $D_j$ is near height zero.  The particular slanted
lower edge is useful: the lower inequality for $D_{j+1}$, after
composition with $J_j$, is precisely $Y>L_j-s$.  Thus $S_j$ lies in
$D_j$, and its translate lies in $D_{j+1}$.  The other two inequalities
have a fixed buffer because $L_j\geq1$ and $s<1/8$.

If $s$ and $\rho$ are each decreased by $\delta$, then the smaller
$D_j$ admits coordinate Cauchy discs of radius at least
$c\delta/(1+A)$ in the larger $D_j$.  The same assertion holds for
seams and for their translated copies.  This follows directly from
\[
 |a_j(p')-a_j(p)|\leq A|p'-p|.
\]
In particular, all fixed-order Cauchy bounds cost only a fixed power
of $\delta^{-1}$, independently of $j$ and $L_j$.  In the estimates
below constants may depend on $h,r,A$ and on fixed fractions of the
initial margins; we take $0<\delta<1$.

\begin{lemma}[Removal of the transverse zero mode]\label{sew:linearize}
Suppose $E=(w+u,p+v)$ is exact symplectic on a seam and
$\|(u,v)\|\leq e$.  After a margin loss $\delta$ there are a
holomorphic function $u_0(p)$ and a periodic mean-zero function $H$
such that
\[
 (u,v)=(u_0,0)+X_H+R,\qquad
 X_H=(H_p,-H_w),\qquad
 \|H\|\leq Ce,\quad \|R\|\leq C\delta^{-2}e^2.
\]
Here $u_0$ is the zero Fourier coefficient of $u$.
\end{lemma}

\begin{proof}
Symplecticity and exactness respectively give
\begin{align}
 u_w+v_p&=-u_wv_p+u_pv_w,\label{sew:symplectic-error}\\
 0&=\int_0^1 v(w,p)(1+u_w(w,p))\,dw.
                                      \label{sew:period-error}
\end{align}
The second identity follows by expanding
$(p+v)d(w+u)-p\,dw$ on a horizontal period circle; the integral of
$p u_w$ is zero.  Consequently its mean $v_0(p)$ satisfies
\[
                  |v_0(p)|\leq C\delta^{-1}e^2.
\]
Take the unique mean-zero periodic primitive satisfying
\[
                        H_w=-v+v_0.
\]
Integration along one horizontal period, followed by subtraction of
the mean, bounds $H$ by $Ce$.  This primitive is holomorphic in $p$.
Equation~\eqref{sew:symplectic-error} gives
\[
 \partial_w(u-H_p)=-u_wv_p+u_pv_w-v_0'(p).
\]
The right side has zero mean and is $O(\delta^{-2}e^2)$ after one
additional fixed fraction of the margin loss.  Integrating its
mean-zero primitive bounds $(u-H_p)-u_0$ by the same quantity.
The other component of $R$ is $v_0$.  Rescaling the allocation of
$\delta$ among these finitely many steps proves the stated estimate.
\end{proof}

The role of exactness is now visible.  Without
Equation~\eqref{sew:period-error}, the first-order constant term in
$v$ need not vanish.  It cannot be absorbed into a translation that
preserves the common transverse coordinate.

\subsection{A uniformly bounded Fourier splitting}

\begin{lemma}[Splitting along the whole chain]\label{sew:split}
Suppose the periodic functions $H_j$ have mean zero, are holomorphic
on $S_j(s,\rho)$, and have norm at most $e$.  Assume
\[
                 |\operatorname{Im}a_j+L_j|\leq\kappa<1/4.
\]
There are periodic holomorphic functions $G_j$ on
$D_j(s-\delta,\rho)$ satisfying
\begin{equation}\label{sew:split-equation}
                  G_j-G_{j+1}\circ J_j=-H_j.
\end{equation}
on the smaller seams.  Their bounds, on domains with the appropriate
additional losses, are
\[
 \|G_j\|\leq C\delta^{-1}e,\qquad
 \|X_{G_j}\|\leq C\delta^{-2}e,\qquad
 \|DX_{G_j}\|\leq C\delta^{-3}e.
\]
All constants are uniform in the block lengths and the number of
blocks.
\end{lemma}

\begin{proof}
Set $T_j=-H_j$ and write
\[
             T_j(w,p)=\sum_{n\neq0}t_{j,n}(p)e^{2\pi i n w}.
\]
The coefficients are holomorphic in $p$.  Denote the negative and
positive parts by $T_j^-$ and $T_j^+$.  Starting from $(w_j,p)$,
continue the coordinates by
\[
                          w_{s+1}=w_s+a_s(p)
\]
in either direction, and define
\begin{equation}\label{sew:fourier-sum}
 G_j(w_j,p)=\sum_{l\geq j}T_l^-(w_l,p)
                         -\sum_{l<j}T_l^+(w_l,p).
\end{equation}
For negative coefficients use a horizontal coefficient-bound line
at height $L_l+s-\delta/4$.  An evaluation point of
$D_j(s-\delta,\rho)$ is below that line, in the continued $l$-th
coordinate, by at least
\[
                     3\delta/4+(1-\kappa)(l-j)
                     \qquad(l\geq j).
\]
For positive coefficients use height $L_l-s+\delta/4$.  The
evaluation point is above that line by at least
\[
                     3\delta/4+(1-\kappa)(j-1-l)
                     \qquad(l<j).
\]
For the nearest terms these inequalities are exactly the two defining
inequalities for $D_j$.  For each further term an intervening block
contributes at least $1-\kappa$.

The absolute value of a Fourier mode is
$|e^{2\pi i n w}|=e^{-2\pi nY}$.  Therefore both sums, including the
sum over frequencies, are bounded by a constant times
\[
 e\sum_{n\geq1}\sum_{q\geq0}
       \exp\bigl(-2\pi n(3\delta/4+(1-\kappa)q)\bigr)
 \leq C e\delta^{-1}.
\]
This proves normal convergence and holomorphicity, including in the
parameter.  Telescoping Equation~\eqref{sew:fourier-sum} leaves
$T_j^-+T_j^+=-H_j$, proving
Equation~\eqref{sew:split-equation}.  Cauchy estimates on smaller
$D_j$ give the derivative bounds.  In particular, one need not
differentiate a sum of an unbounded number of translation functions
term by term.
\end{proof}

\subsection{The nonlinear iteration}

We record the domain issues in the Newton step.  Reserve a fixed
number of successive margin losses, each comparable to $\delta$,
for the primitive, Fourier splitting, derivative estimates, flows,
and composition.  Enlarging this fixed number if necessary does not
depend on any block length.  In the following statements the total
loss is denoted by $c_0\delta$, with $c_0$ a fixed constant.

Apply Lemma~\ref{sew:linearize} to $E_j=J_j^{-1}F_j$, and then
Lemma~\ref{sew:split}.  Let $K_j$ be the time-one map of $X_{G_j}$.
It is defined on the reserved smaller $D_j$ provided
$C\delta^{-2}e<\delta$.  Its displacement and first-order remainder
satisfy
\begin{equation}\label{sew:flow-estimates}
 \|K_j-\operatorname{id}\|\leq C\delta^{-2}e,\qquad
 \|K_j-\operatorname{id}-X_{G_j}\|
                              \leq C\delta^{-5}e^2.
\end{equation}
Indeed the second estimate follows by integrating the differential
equation for the flow and using
$\|DX_{G_j}\|\,\|X_{G_j}\|\leq C\delta^{-5}e^2$.
Negative-time flows satisfy the same estimates on the corresponding
smaller domains.

The new transition and its proposed model are
\[
 F_j^{\mathrm{new}}=K_{j+1}^{-1}F_jK_j,\qquad
 J_j^{\mathrm{new}}(w,p)=(w+a_j(p)+u_{j,0}(p),p).
\]
To check the cancellation without differentiating a long coordinate
change, put $Q_j=G_{j+1}\circ J_j$.  The conjugate
$J_j^{-1}K_{j+1}J_j$ is the time-one Hamiltonian flow of $Q_j$.
On a buffered seam $\|Q_j\|\leq C\delta^{-1}e$; the same Cauchy
and flow estimates apply to it.  The linear term of
\[
        J_j^{-1}K_{j+1}^{-1}J_j\, E_j\, K_j
\]
is
\[
 (u_j,v_j)+X_{G_j}-X_{Q_j}
            =(u_{j,0},0)+R_j.
\]
Here the Hamiltonian terms cancel by
Equation~\eqref{sew:split-equation}.  The remaining products are
quadratic: Equation~\eqref{sew:flow-estimates} bounds the flow
remainders by $C\delta^{-5}e^2$, and the other composition errors
are bounded by a first derivative times a displacement.  For example,
they are bounded by
\[
 C(\delta^{-1}e)(\delta^{-2}e)
       \quad\hbox{or}\quad
 C(\delta^{-3}e)(\delta^{-2}e).
\]
Finally compose with $(w,p)\mapsto(w-u_{j,0}(p),p)$ to normalize
by $J_j^{\mathrm{new}}$.  Its derivative differs from the identity
by at most $C\delta^{-1}e$ and introduces only another quadratic
term.  It follows, for some fixed integer $q\geq6$, that
\begin{equation}\label{sew:newton-bound}
 \sup_j\|(J_j^{\mathrm{new}})^{-1}F_j^{\mathrm{new}}
                                     -\operatorname{id}\|
                              \leq C\delta^{-q}e^2.
\end{equation}
One can use the displayed estimates with $q=6$; allowing a larger
fixed $q$ absorbs harmless reallocations of the margins.  No
operation in this step involves more than two adjacent charts.

The new left boundary of $D_j$ changes by at most
$\|u_{j-1,0}\|\leq e$.  Flow displacements are
$O(\delta^{-2}e)$.  Hence all new domains, translated seams and
inverse-flow domains fit in their reserved predecessors when
$C\delta^{-q}e$ is sufficiently small compared with $\delta$.
The model derivative changes by at most $C\delta^{-1}e$.
These observations justify both the domain inclusions used above and
their uniform repetition.

Exactness persists.  With $\lambda=p\,dw$, Cartan's formula for our
Hamiltonian convention gives
\[
                  \mathcal L_{X_H}\lambda=d(pH_p-H).
\]
Thus the flows are exact.  The shear $J_j$ is exact because
$J_j^*\lambda-\lambda=p a_j'(p)\,dp$ has a holomorphic primitive
on the parameter disc.  Compositions and inverses of exact maps are
exact.

\begin{proof}[Proof of Theorem~\ref{sew:chain}]
Start with slightly smaller $s_0<h$ and $\rho_0<r$.  Fix a small
$d_0>0$ and take $\delta_n=d_0 2^{-n}$.  Choose $d_0$ so that
\[
 c_0\sum_{n\geq0}\delta_n
                   <\tfrac14\min(s_0,\rho_0).
\]
At stage $n$ reduce both margins by $c_0\delta_n$.  Thus their
limits remain positive.  Set $B=C d_0^{-q}$, increasing the uniform
constant in Equation~\eqref{sew:newton-bound} once if needed.  If
$e_n$ bounds the normalized errors, then
\[
 e_{n+1}\leq B2^{qn}e_n^2.
\]
For $x_n=B2^{q(n+1)}e_n$ this becomes $x_{n+1}\leq x_n^2$.
Consequently, if $x_0\leq\theta<1$, then
\[
 e_n\leq B^{-1}2^{-q(n+1)}\theta^{2^n}.
\]
Every series $\sum_n\delta_n^{-k}e_n$, with fixed $k$, converges
and tends to zero as $e_0\to0$.  Taking $e_0$ still smaller ensures
all the margin conditions above, keeps
$|\operatorname{Im}a_j^{(n)}+L_j|<h/3$, and keeps the translation
derivative bound below $A+1$.  These estimates are uniform in $j$.

The successive compositions of the $K_j$ converge uniformly, with
derivatives on smaller domains, to maps $C_j$ into the original
block domains.  To see this directly, sum their displacements;
their derivative norms are bounded by a convergent product of
$1+C\delta_n^{-3}e_n$.  The reserved inclusions keep every
composition in its domain.  The limits remain symplectic because
their differentials converge and each determinant is one.  Put
$\widetilde\phi_j=\phi_j\circ C_j$ and
$\widetilde a_j=\lim_n a_j^{(n)}$.  Passing to the limit in the
chart identities on the seams gives
\[
 \widetilde\phi_j(w,p)=
       \widetilde\phi_{j+1}(w+\widetilde a_j(p),p).
\]
The parameter radius has been decreased by the fixed sum of reserved
losses and is still bounded below by a positive number depending only
on the initial margins. A smaller common radius $b$ can therefore be
fixed before the error threshold, independently of the chain.

It remains to identify the domain, rather than invoke a
uniformization theorem.  Fix block zero, put $w_0=w$, and define
$w_{j+1}=w_j+\widetilde a_j(p)$ in both directions.  These are
holomorphic shears, and $dw_j\wedge dp=dw\wedge dp$.  Use the
$j$-th chart between the heights
\[
       \operatorname{Im}w_{j-1}=L_{j-1},\qquad
       \operatorname{Im}w_j=L_j,
\]
including the surviving overlaps.  In $w_j$ coordinates the first
height is near zero and the second is $L_j$.  Their separation is
at least $1-h/3$.  Hence the successive edges tend to both
infinities, locally uniformly in $p$, and the charts cover all of
$\C\times\Delta_b$.  They agree on overlaps by the limiting
identity.  Their common pullback form is $dw\wedge dp$, proving
immersion.  Periodicity follows from the periodicity of the original
charts and every correction.  The summable estimates also prove
the claimed approximation statements.
\end{proof}

\subsection{Enforcing exactness by recentering}

The geometric applications supply nearly matching blocks, but not
initially exact transitions.  The following observation explains
precisely how their centers are allowed to move.

\begin{lemma}[Period adjustment]\label{sew:centers}
Suppose symplectically normalized charts, together with corresponding
model charts, are available for centers varying in larger transverse
discs.  Suppose the actual charts are uniformly close to their
models, and transitions at matching model positions are uniformly
close to degree-one shears.  The transition into a newly added block
can be made exact by a transverse translation of that block of size
bounded by its transition error.  Recenter its model by the same
translation.  The approximation of that block by its own recentered
model retains its original bound.

Consequently, such adjustments can be repeated indefinitely whenever
the updated centers remain in a region with uniform buffered chart
choices.  No bound on their accumulated displacement relative to
the first center is required.
\end{lemma}

\begin{proof}
For a transition $F=(W,P)$ the closed one-form
$P\,dW-p\,dw$ has a constant complex period
\[
                     D=\int(P\,dW-p\,dw).
\]
It is independent of the representative period circle and of $p$.
Evaluation at an interior circle with $p=0$, using Cauchy estimates
on a fixed seam margin, bounds $|D|$ by the error from a shear.
Replace the added target chart by
$(W,P_{\rm new})\mapsto\phi(W,P_{\rm new}+D)$.  The transition
now has $P_{\rm new}=P-D$, so its period is
\[
        D-D\int dW=0.
\]
The last integral is one because the transition has degree one.
In a backward extension one instead translates the newly added
input chart, with the corresponding sign.  Uniform buffers permit
these small translations.

Apply the identical substitution to the model chart.  Its outgoing
end is now the outgoing end of that updated model, and this is the
position used to choose the following block.  Both actual and
model maps were defined on a larger disc, so their difference after
restriction and translation has the same uniform bound.  The old
seam still has only a small error; the assertion does not require
the two moved model centers to coincide there.  Repeating this
argument compares each new chart with its current model and never
sums the previous matching errors.
\end{proof}

For completeness, the normalization used in this lemma is compatible
with periodicity.  If a nearby chart pulls back the surface form to
$k(w,t)\,dw\wedge dt$, with $k$ periodic and close to one, put
\[
                       p(w,t)=\int_0^t k(w,s)\,ds.
\]
Then $dw\wedge dp=k\,dw\wedge dt$, and $(w,t)\mapsto(w,p)$ is a
small periodic biholomorphic change on smaller domains.  Its inverse
is obtained uniformly with margins.  Additive origins in $w$ may
also be changed by holomorphic functions of $p$; these changes are
exact shears, so they do not affect the period adjustment.
